% !PRINTED 605
\weberpage{605}

Darin kann $\chi$ jeder der Charaktere $\chi_\kappa$ sein, und die Summe erstreckt sich über sämtliche Ideale $\mathfrak{a}$ von $\overline{O}$. Nach (11) ist, wenn wir

\[
G_\kappa = \sum^{i}\chi_\kappa(\overline{A}_i)\,C_i
\]

setzen,

\weq{17}{\begin{aligned}
Q_1 &= \frac{hG}{s-1} + G_1, \\
Q_\kappa &= G_\kappa, \quad \kappa > 1.
\end{aligned}}
% [source sic: equation (17) prints hG, whereas equation (21) later uses hg for the same residue]

worin die $G_1, G_2, \dots, G_h$ \gesperrt{Funktionen von $s$ sind, die für $s = 1$ einen endlichen Wert behalten.}

Die Summen $Q$ lassen sich in derselben Weise umformen, die wir schon im \S\,197 des zweiten Bandes kennen gelernt haben.

Danach ist, wenn $\mathfrak{p}$ irgend ein Primideal in $\overline{O}$ bedeutet:

\[
\begin{aligned}
\frac{1}{1 - \chi(\mathfrak{p})N(\mathfrak{p})^{-s}} &= 1 + \chi(\mathfrak{p})N(\mathfrak{p})^{-s} + \chi(\mathfrak{p})^2 N(\mathfrak{p})^{-2s} + \cdots \\
&= 1 + \frac{\chi(\mathfrak{p})}{N(\mathfrak{p})^s} + \frac{\chi(\mathfrak{p}^2)}{N(\mathfrak{p}^2)^s} + \frac{\chi(\mathfrak{p}^3)}{N(\mathfrak{p}^3)^s} + \cdots,
\end{aligned}
\]

und da sich alle Ideale $\mathfrak{a}$ als Produkte solcher Ideale $\mathfrak{p}$ darstellen lassen, so folgt:

\weq{18}{Q = \Pi\,\frac{1}{1 - \chi(\mathfrak{p})N(\mathfrak{p})^{-s}},}

worin sich das unendliche Produkt $\Pi$ auf alle Primideale $\mathfrak{p}$ in $\overline{O}$ erstreckt.

Zu jedem Ideal $\mathfrak{a}$ in $\overline{O}$ gehört ein gewisser Exponent $\varphi$, d.\,h.\ es gibt einen gewissen niedrigsten positiven Exponenten $\varphi$, für den $\mathfrak{a}^\varphi$ in der Hauptklasse $\overline{A}_1$ enthalten ist. Die Zahl $\varphi$ ist immer ein Teiler der Klassenzahl $h$. Wir setzen:

\weq{19}{h = \varepsilon\varphi.}

Gehört in diesem Sinne $\mathfrak{a}$ zum Exponenten $\varphi$, so sind sämtliche Charaktere $\chi_k(\mathfrak{a})$ $\varphi$te Einheitswurzeln, und jede $\varphi$te Einheitswurzel kommt darunter $\varepsilon$ mal vor (Bd.\ II, \S\,198, 1.).
% [source sic: the index switches here from \chi_\kappa to \chi_k]

Gehört also $\mathfrak{p}$ zum Exponenten $\varphi$, so ist

\[
\begin{aligned}
&[1 - \chi_1(\mathfrak{p})N(\mathfrak{p})^{-s}][1 - \chi_2(\mathfrak{p})N(\mathfrak{p})^{-s}]\cdots[1 - \chi_h(\mathfrak{p})N(\mathfrak{p})^{-s}] \\
&\qquad = [1 - N(\mathfrak{p})^{-\varphi s}]^\varepsilon,
\end{aligned}
\]

und es folgt aus (18):

\weq{20}{Q_1 Q_2 \cdots Q_h = \Pi\,\frac{1}{[1 - N(\mathfrak{p})^{-\varphi s}]^\varepsilon}.}
